\begin{abstract}
Benchmark audits are often interpreted as if their diagnostics were already validated measurement
instruments. We study a stricter alternative: each diagnostic must expose its sensitivity,
uncertainty, practical materiality, external validity, and transfer limits before it can license a
benchmark-quality or model-ranking claim. We implement this discipline in ValidEval through
versioned evidence states, exact-identity contracts, leakage guards, null-calibrated ranking
analyses, and fail-closed import and reporting paths. As a historical case study, we independently
reconstruct a public HELM-derived MMLU response panel with 39 models, 14,042 items, and no missing
cells. The reconstruction verifies descriptive rank variation but does not treat a legacy raw rank
threshold as evidence of material instability. Controlled cross-benchmark, human, external-label,
and decoupled-synthetic results remain explicit result requirements. The intended contribution is
therefore a falsifiable measurement-audit methodology, not a scalar benchmark-validity score.
\end{abstract}
